\section{Experiments and Results}
% Source of the artifacts shown in this figure (winning harness per dataset):
%   TB2:    runs/exp-rho-tb2/harness/h_edccd200e74a/
%   SWE:    runs/exp-rho-swebench/harness/h_c13b7933bad0/
%   GAIA-2: runs/exp-rho-gaia2-patched/harness/h_abff8c30ecd6/
% See docs/paper/figures/fig-harness-artifacts.html for the figure source and
% inline mapping from each artifact item back to its README/script.
\begin{figure*}[!t]
\centering
\includegraphics[width=\textwidth]{figures/fig-harness-artifacts.pdf}
\caption{\looseness=-1 Highest-scoring harness produced by \method{} on each benchmark. \emph{Instructions} are task-agnostic procedural rules, \emph{Skills} record grader or environment idiosyncrasies that previously caused failures, and \emph{Tools} are executable scripts. Items shown are representative, and the full verbatim contents of each harness are in Appendix~\ref{app:artifacts}.}
\label{fig:harness-artifacts}
\vspace{-1em}
\end{figure*}



\noindent\textbf{Setup.}
We use the Codex agent \citep{openai2025codex} as the base harness for retrospective optimization.
Specifically, this agent uses GPT-5.5 \citep{openai2026gpt55} configured with high reasoning effort.
When we invoke Codex to solve a task, we construct the harness as a configurable workspace folder.
This folder contains executable scripts as tools, along with text files for skills and instructions.
For all our experiments, we set the coreset size $k$ to 10.
In addition, we use 3 for both parallel trajectory sampling and harness proposals.
To measure the improvement, we report the pass rate on the held-out test set using both the vanilla Codex harness and the optimized one.

\noindent\textbf{Data.}
We collect past trajectories from existing benchmark datasets.
Specifically, we divide the original benchmarks into a trajectory set and a test set.
We then run the vanilla Codex agent on the trajectory set to generate the required trajectories for \method{}.
Building on this, we evaluate \method{} on SWE-Bench Pro, Terminal-Bench~2, and \mbox{GAIA-2}.
SWE-Bench Pro contains long-horizon software-engineering tasks requiring repository-level reasoning and multi-file edits \citep{deng2025swebenchpro}.
Terminal-Bench~2 contains command-line tasks with executable graders \citep{terminalbench2}.
\mbox{GAIA-2} evaluates LLM agents in dynamic, asynchronous environments for knowledge work \citep{froger2026gaia2}.
As a result, these three benchmarks cover a wide range of task types across software engineering, technical work, and knowledge work.
During optimization \method{} consumes only past trajectories and its own judgments; graders are used solely for held-out evaluation.
We provide detailed information about the benchmarks and data splits in Appendix~\ref{app:datasets}.

\subsection{Comparison with Feedback-Free Baselines}

We compare \method{} against three competitive harness optimization methods that do not require validation feedback.
For the baselines, \textit{Dynamic Cheatsheet} maintains a running record of useful facts and procedures \citep{suzgun2025dynamiccheatsheet}.
\textit{ReasoningBank} stores reusable reasoning patterns and retrieves the top-$k$ relevant entries at inference time \citep{ouyang2026reasoningbank}.
Similarly, \textit{Sleep-time Compute} preprocesses past traces offline into compact notes, which are then prepended to the agent's context \citep{lin2025sleeptime}.
We adapt each baseline to our datasets and agent setting while holding the total agent-call budget approximately fixed to ensure a fair comparison.
Detailed adaptation procedures are provided in Appendix~\ref{app:baselines}.

As Table~\ref{tab:main} shows, \method{} delivers consistent improvements across all three benchmarks, whereas the baselines do not.
Most notably, we achieve an absolute improvement of 19\% on SWE-Bench Pro without relying on any validation-based grading.
We attribute this advantage to the more flexible harness optimization that \method{} enables.
Specifically, the agent can create new tools, skills, and instructions for the harness, whereas previous methods focus predominantly on memory systems or text-based skills.
Furthermore, the use of self-preference may contribute to the consistency of these harness improvements.
In contrast, the performance gains of the baseline methods tend to be smaller and vary across different datasets.
In the next section (Section~\ref{sec:case-study}), we examine how \method{} modifies the harness to improve the agent.




\subsection{What the Optimized Harness Contains}
\label{sec:case-study}

Figure~\ref{fig:harness-artifacts} summarizes and interprets the new harness contents generated after \method{} optimization.
In our work, the harness is materialized as a directory containing markdown files for instructions and skills, as well as executable scripts for tools.

Across all three benchmarks, \method{} adds multiple new skills and tools to the harness.
Many of these new additions address typical failure modes encountered by the original harness.
For example, in SWE-Bench Pro, the agent learns that the Go toolchain resides at a non-standard location outside the default path.
It also discovers that Python cache directories must be stripped before producing the final diff, as failing to do so often prevents patches from applying cleanly.
To address these, the agent adds a new \texttt{check\_build\_and\_lint} tool that locates non-standard toolchains and flags the generated artifacts that must be kept out of the patch, fixing the diff-hygiene procedures the original trajectories repeatedly missed.
These examples illustrate how \method{} identifies useful tools and skills across diverse scenarios by analyzing past failures.

\subsection{Comparison with Validation-Feedback Optimization}

We next compare \method{} against Meta-Harness \citep{lee2026metaharness}.
Meta-Harness is a validation-feedback optimizer: a proposer emits harness edits, each candidate is graded on a labeled search set with the ground-truth grader, and the search keeps the harnesses that score best on that set.
To maintain a fair comparison, we use the same Codex agent as the Meta-Harness proposer and solver.
Unlike \method{}, Meta-Harness requires held-out labels, and because it evolves over multiple rounds, it consumes more agent calls.
Therefore, we evaluate it at a single round to match our compute budget, as well as in an extended-budget setting of ten rounds.

\begin{table}[t]
\caption{\method{} versus Meta-Harness, a validation-feedback optimizer, on SWE-Bench Pro. \emph{Agent calls} is the optimization-time budget relative to one \method{} run. Held-out pass rate on the same split as Table~\ref{tab:main}.}
\label{tab:meta-harness}
\centering
\footnotesize
\setlength{\tabcolsep}{4pt}
\resizebox{\columnwidth}{!}{%
\begin{tabular}{lccc}
\toprule
\textbf{Method} & \shortstack{\textbf{Val.}\\\textbf{labels}} & \shortstack{\textbf{Agent}\\\textbf{calls}} & \shortstack{\textbf{SWE-Bench}\\\textbf{Pro}} \\
\midrule
\rowcolor{rhoBlue}
\method{} & none & 103 ($1.0\times$) & 0.78 \\
\rowcolor{white}
Meta-Harness (1 round) & required & 41 ($0.4\times$) & 0.62 \\
\rowcolor{rowAlt}
Meta-Harness (10 rounds) & required & 320 ($3.1\times$) & \textbf{0.80} \\
\bottomrule
\end{tabular}%
}
\vspace{-1.5em}
\end{table}

At the matched single-round budget, Meta-Harness selects its best candidate using validation scores but achieves only a 0.62 pass rate on SWE-Bench Pro.
This is substantially lower than the 0.78 pass rate achieved by \method{}.
Scaling Meta-Harness to a 10-round setting increases its performance ceiling to 0.80 on SWE-Bench Pro.
However, this higher performance requires roughly three times the optimization-phase compute of \method{} (320 vs.\ 103 agent calls) and, more importantly, still depends on held-out labels that \method{} does not use.
